\section{Tag 2 — Cleverere Prüfziffern und der Sprung zum Korrigieren}

\subsection*{Bleistiftübung 1 — ISBN-10}

(b) \texttt{0306406152}:
\begin{verbatim}
0·10 + 3·9 + 0·8 + 6·7 + 4·6 + 0·5 + 6·4 + 1·3 + 5·2
= 0 + 27 + 0 + 42 + 24 + 0 + 24 + 3 + 10 = 130
130 mod 11 = 9    (denn 130 = 11·11 + 9)
11 − 9 = 2
\end{verbatim}
Prüfziffer = \textbf{2}, also \texttt{0-306-40615-2}. Das ist „Gödel,
Escher, Bach“ von Douglas Hofstadter — ein Klassiker.

(c) \texttt{097522980}:
\begin{verbatim}
0·10 + 9·9 + 7·8 + 5·7 + 2·6 + 2·5 + 9·4 + 8·3 + 0·2
= 0 + 81 + 56 + 35 + 12 + 10 + 36 + 24 + 0 = 254
254 mod 11 = 1    (denn 254 = 23·11 + 1)
11 − 1 = 10
\end{verbatim}
Prüfziffer = \textbf{10}, geschrieben als \textbf{X}. Vollständige
ISBN: \texttt{097522980X}.

\subsection*{Bleistiftübung 2 — Warum Modulo 11}

(a) Differenz: $(b - a)(11 - i) + (a - b)(11 - j) = (a - b)((11-j) -
(11-i)) = (a - b)(i - j)$.

(b) Damit der Tausch unentdeckt bleibt, müsste $(a - b)(i - j) \equiv
0 \pmod{11}$. Da 11 \textbf{prim} ist, gilt: ein Produkt ist nur dann
durch 11 teilbar, wenn einer der Faktoren durch 11 teilbar ist. Aber
$a - b$ liegt zwischen $-9$ und $9$ (und ist nicht 0, weil $a \ne b$),
und $i - j$ liegt zwischen $-8$ und $8$ (und ist nicht 0, weil
verschiedene Positionen). Beide Faktoren sind also \textbf{niemals
durch 11 teilbar} — der Tausch wird immer erkannt!

(c) Bei Modulo 10 funktioniert das Argument nicht, weil $10 = 2 \cdot
5$ (keine Primzahl). Beispiel: $a-b = 5$ und $i-j = 2$ ergibt 10 →
Tausch unentdeckt. Genau das Phänomen aus Tag 1.

\subsection*{Bleistiftübung 3 — Detektivarbeit}

ISBN: \texttt{0-30?-40615-2} mit unbekannter Ziffer $x$ an Position 4.

Gewichtete Summe inklusive Prüfziffer 2:
\begin{verbatim}
0·10 + 3·9 + 0·8 + x·7 + 4·6 + 0·5 + 6·4 + 1·3 + 5·2 + 2·1
= 0 + 27 + 0 + 7x + 24 + 0 + 24 + 3 + 10 + 2
= 90 + 7x
\end{verbatim}
Damit die ISBN gültig ist, muss diese Summe durch 11 teilbar sein:
\begin{align*}
  90 + 7x &\equiv 0 \pmod{11} \\
  90 &\equiv 2 \pmod{11} \\
  \Rightarrow 2 + 7x &\equiv 0 \pmod{11} \\
  7x &\equiv -2 \equiv 9 \pmod{11}
\end{align*}
Probiere $x = 0, 1, \dots, 9$: $x = 6$ ergibt $7 \cdot 6 = 42 \equiv 9
\pmod{11}$ \checkmark.

Die fehlende Ziffer ist \textbf{6}, die ISBN lautet
\texttt{0-306-40615-2}.

\subsection*{Aufgabe 1.1 — Lösung}

\pythoncodeteil{tag2/isbn10.py}{zifferndreher}

Erwartung: keine unentdeckten Vertauschungen.

\subsection*{Bleistiftübung 4 — Luhn}

(a) \texttt{4716 7300 5852 1908}:
Von rechts nach links Position-Index 0, 1, 2, …
\begin{verbatim}
8 (× nicht)
0 → 0·2 = 0
9 (× nicht)
1 → 1·2 = 2
2 (× nicht)
5 → 5·2 = 10 → 1
8 (× nicht)
5 → 5·2 = 10 → 1
0 (× nicht)
0 → 0·2 = 0
3 (× nicht)
7 → 7·2 = 14 → 5
6 (× nicht)
1 → 1·2 = 2
7 (× nicht)
4 → 4·2 = 8

Summe = 8+0+9+2+2+1+8+1+0+0+3+5+6+2+7+8 = 62
62 mod 10 = 2 ≠ 0 → ungültig
\end{verbatim}
Die Nummer ist \textbf{nicht} gültig.

(c) Die meisten Einzelfehler werden erkannt. Tatsächlich: Luhn erkennt
\textbf{alle} Einzelfehler in einer einzelnen Ziffer.

\subsection*{Aufgabe 2.1 — Welche Tippfehler bleiben unentdeckt?}

\pythoncodeteil{tag2/luhn.py}{zifferndreher}

Ergebnis: die einzige Vertauschung benachbarter Ziffern, die Luhn
\textbf{nicht} erkennt, ist $09 \leftrightarrow 90$. Alle anderen
werden erkannt.

\subsection*{Bleistiftübung 5 — Distanzen}

\begin{center}
\begin{tabular}{llc}
  \toprule
  A & B & Distanz \\
  \midrule
  \texttt{1011010}  & \texttt{1011010}  & 0 \\
  \texttt{0000}     & \texttt{1111}     & 4 \\
  \texttt{11001100} & \texttt{10101010} & 4 \\
  HUND              & RUND              & 1 \\
  TANNE             & KANNE             & 1 \\
  \texttt{0000000}  & \texttt{1111111}  & 7 \\
  \bottomrule
\end{tabular}
\end{center}

\subsection*{Bleistiftübung 6 — Mindestdistanz}

(a) Schauen wir z.\,B.\ \texttt{0000} vs.\ \texttt{0011}: Distanz 2.
Oder \texttt{0011} vs.\ \texttt{0101}: Distanz 2. Tatsächlich ist die
Mindestdistanz beim Paritätscode immer \textbf{2} (jedes Codewort
unterscheidet sich vom nächstgelegenen genau in 2 Bits, weil jede
Veränderung von 1 Bit die Parität kippen würde).

(b) Wiederholungscode $(000, 111)$ hat Mindestdistanz \textbf{3}, also
größer als der Paritätscode (2).

\subsection*{Bleistiftübung 7 — Tabelle}

\begin{center}
\begin{tabular}{lccc}
  \toprule
  Code & Mindestdistanz & Erkennbar & Korrigierbar \\
  \midrule
  Paritätscode $(3+1)$        & 2 & 1 & 0 \\
  Wiederholungscode $(3,1)$   & 3 & 2 & 1 \\
  Wiederholungscode $(5,1)$   & 5 & 4 & 2 \\
  Wiederholungscode $(7,1)$   & 7 & 6 & 3 \\
  \bottomrule
\end{tabular}
\end{center}

Beobachtung: bei Wiederholungscodes der Länge $n$ ($n$ ungerade) ist
die Mindestdistanz $n$, also können $(n-1)/2$ Fehler korrigiert
werden. Sehr robust, aber sehr ineffizient.

\subsection*{Aufgabe 3.1 — Lösung}

\pythoncodeteil{tag2/hamming_distanz.py}{mindestdistanz}

\subsection*{Aufgabe 3.3 — Doppelfehler beim Wiederholungscode}

\pythoncodeteil{tag2/hamming_distanz.py}{doppelfehler}

Ergebnis: 0/6 (0\,\%!). Bei 2 Fehlern liegt das Wort näher am falschen
Codewort. Das ist genau, was die Theorie vorhersagt: Distanz 3 →
korrigierbar nur 1 Fehler.

\subsection*{Cliffhanger-Knobelaufgabe}

Bei 4 Datenbits gibt es $2^4 = 16$ mögliche Datenwörter. Jedes davon,
kodiert in $n$ Bits, hat $n$ mögliche 1-Bit-Nachbarn. Damit alle 16
Codewörter und ihre $n$ Nachbarn nicht überlappen, brauchen wir
mindestens $16 \cdot (n+1) \le 2^n$.

\begin{itemize}
\item $n = 5$: $16 \cdot 6 = 96$, aber $2^5 = 32$ → reicht nicht.
\item $n = 6$: $16 \cdot 7 = 112$, aber $2^6 = 64$ → reicht nicht.
\item $n = 7$: $16 \cdot 8 = 128$, und $2^7 = 128$ → \textbf{passt
  genau}.
\end{itemize}

Antwort: \textbf{7 Bits} (also 3 Prüfbits zusätzlich zu den 4
Datenbits). Das ist der berühmte \textbf{(7,4)-Hamming-Code} — und das
Erstaunliche: die Schranke wird \emph{exakt} getroffen, kein Bit ist
zu viel.
